\documentclass[aps,prl,onecolumn,superscriptaddress]{revtex4-2}
\usepackage{amsmath,amssymb,amsfonts}
\usepackage{booktabs}
\usepackage{array}
\usepackage{xcolor}
\usepackage{hyperref}

\definecolor{navy}{HTML}{0B2545}
\definecolor{dgrn}{HTML}{1B5E20}
\definecolor{teal}{HTML}{006064}
\definecolor{purp}{HTML}{4527A0}
\definecolor{mgrey}{HTML}{555555}

\begin{document}

\title{Appendix KC1: The BCT Periodic Table ---
Topological Classification of All 118 Elements\\
\normalsize Companion to BCT Letter 84}

\author{Michel Robert Cabri\'{e}}
\affiliation{Independent Researcher, Barrys Reef, Victoria, Australia}
\affiliation{ORCID: 0009-0007-9561-9859}

\date{20 March 2026}

\begin{abstract}
We provide the complete topological classification of all 118 elements
under the BCT Superfluid Lattice Model, organised by Hopf multiplet
structure. Eleven classes are defined, each with a rigorous BCT
derivation. We justify the Janet (left-step) layout as the natural
periodic arrangement consistent with OHC Bessel resonance mode ordering,
and define the Hopf fill fraction that annotates each element in the
BCT Periodic Table figures.
\end{abstract}

\maketitle

\section{Classification scheme}

Each element is assigned to one of eleven BCT topological classes
based on the Hopf multiplet structure of its outermost active modes.
The classification follows directly from the OHC mode theorems
(Appendices KA1, KA4) and the Hopf symmetry theorems (KB1, KB2, KB3).

\begin{table}[h]
\centering
\caption{BCT topological classification of the elements.
$^\dagger$ Carbon appears in both Life axis and p-block;
Life axis takes precedence.
$^\ddagger$ Lu and Lr appear in f-closure; f-block
refers to remaining lanthanides and actinides.}
\begin{tabular}{p{26mm}p{58mm}p{50mm}}
\toprule
Class & Defining BCT property & Representative elements \\
\midrule
Noble gas    & Closed Hopf multiplet. No vacant modes. Bonding topologically impossible (L81). & He, Ne, Ar, Kr, Xe, Rn, Og \\[4pt]
Life axis    & $S_H = 1$ at sp$^3$ half-fill. Unique equal filling of Hopf channels (L82). & H, C, Si, Ge, Sn, Pb, Fl \\[4pt]
d half-fill  & $S_H(d) = 0.50$ at $N_d = 5$. d-shell Hopf symmetry attractor (L83). & Cr, Mo, Mn, Tc, Re \\[4pt]
d closure    & Closed d Hopf multiplet, $N_d = 10$. d-shell closure attractor (L83). & Cu, Ag, Au, Pd, Zn, Cd, Hg, Cn \\[4pt]
f half-fill  & $S_H(f) = 1.00$ at $N_f = 7$. f-shell Hopf attractor (KB3). & Eu, Am \\[4pt]
f closure    & Closed f Hopf multiplet, $N_f = 14$ (KB3). & Yb, No, Lu, Lr \\[4pt]
Ferromagnetic & Collective Hopf alignment in d-shell (L83). & Fe, Co, Ni \\[4pt]
s-block      & Outermost active mode is $r_\mathrm{oct}$ (octahedral void). & Li, Be, Na, Mg, K, Ca\ldots \\[4pt]
p-block      & Outermost active modes are $r_\mathrm{tet}$ (tetrahedral void). & B, N, O, F, Al\ldots \\[4pt]
d-block      & Active d-modes, not at Hopf symmetry point. & Sc, Ti, V, and others \\[4pt]
f-block      & Active f-modes. & La--Yb$^\ddagger$, Ac--No$^\ddagger$ \\
\bottomrule
\end{tabular}
\end{table}

\section{Priority rule for overlapping classes}

Some elements satisfy criteria for more than one class.
The priority rule is: Noble gas $>$ Life axis $>$ d/f-closure
$>$ d/f-half-fill $>$ Ferromagnetic $>$ d-block $>$ f-block
$>$ p-block $>$ s-block. The higher-priority class is assigned.
Carbon (Life axis) takes priority over p-block; Cu (d-closure)
takes priority over d-block.

\section{The Janet left-step layout}

The standard Mendeleev table places the s-block at the left
and detaches the f-block as a footnote. The Janet (left-step)
layout, first proposed by Charles Janet in 1928, places the
s-block at the right and the f-block at the left.
The reading order left-to-right is: f $\to$ d $\to$ p $\to$ s.

In the BCT framework, this is not cosmetic.
The Janet order corresponds exactly to the OHC Bessel resonance
filling sequence -- from deepest Planck-scale mode (f, innermost
tetrahedral void resonance) to outermost mode (s, principal
octahedral void). The Janet layout is the \emph{natural} BCT layout.

Two additional consequences follow:
\begin{enumerate}
\item Helium sits above Beryllium (both $s^2$ configurations)
  rather than above Neon. This is correct: He is a closed
  1s$^2$ shell (Hopf closure at $n=1$), not a p-shell element.
\item The noble gas column forms the \emph{right edge} of every
  period -- geometrically marking the topological closure of each
  Hopf shell.
\end{enumerate}

\section{The Hopf fill fraction}

Each element cell in the BCT Periodic Table displays a Hopf fill
fraction bar (linear layout) or arc (circular layout), defined as:
%
\begin{equation}
  f_H = \frac{N_\mathrm{active}}{N_\mathrm{capacity}}
\end{equation}
%
where $N_\mathrm{active}$ is the electron count in the outermost
active Hopf shell and $N_\mathrm{capacity}$ is the shell capacity
(2 for s, 6 for p, 10 for d, 14 for f).
This provides an immediate visual indicator of position within
the topological period:

\begin{align}
  f_H = 0.00 &\quad \text{(empty shell, e.g.\ La: }4f^0\text{)} \nonumber\\
  f_H = 0.50 &\quad \text{(half-fill: Life axis, Cr, Eu)} \nonumber\\
  f_H = 1.00 &\quad \text{(closed: noble gases, Cu, Yb)}
\end{align}

Noble gases display a full disc ($f_H = 1.00$);
life axis elements display a half-arc ($f_H = 0.50$);
all other elements display a proportional arc.

\section{Companion figures}

This appendix accompanies two publication figures:
\begin{itemize}
\item \texttt{BCT\_PeriodicTable.pdf} -- Janet left-step layout,
  A3 landscape, all 118 elements colour-coded by BCT class.
\item \texttt{BCT\_PeriodicTable\_Round\_v12.pdf} -- circular
  in-the-round layout, A3, noble gases at 12~o'clock, Life axis
  at 6~o'clock, Hopf fill arcs in every element circle.
\end{itemize}
Both figures are available open-access on Zenodo
(search: BCT superfluid lattice, ORCID 0009-0007-9561-9859).

\end{document}
